% page 275 of the facsimile (image p-274 of the scan)
% block 165.1 x 237.7 mm ; left margin 36.9 mm
\pgc{15.240mm}{0.000mm}{
\l{\cel{55}267}
\l{}
\l{\sou{karnaciono}. Koloro di la parti extera di la homo-}
\l{\cel{3}korpo. - DEFIS.}
\l{}
\l{\sou{karnavalo}. (\sou{religio katol.})\cel{2}La lasta karno-dii di la periodo de}
\l{\cel{3}la Reji-dio til la Cindro-merkurdio, e dum qui eventas la la-}
\l{\cel{3}sta amuzi e distrakti di ta periodo ; travestii, bali de mas-}
\l{\cel{3}kiti, e c. - DEFIRS.}
\l{}
\l{\sou{karnivoro}.\cel{2}Qua su nutras ye karno. - DEFIS.}
\l{}
\l{\sou{karonado}.\cel{2}Pafilo, ordinare ek fero, plu kurta kam kanono, qua}
\l{\cel{3}uzesis, precipue sur la navi, e qua (quale la bombardilo)}
\l{\cel{3}havas alveolo por la pulvero. - DEFIS.}
\l{}
\l{\sou{karoso}.\cel{2}Luxo-veturo, di qua la korpo su apogas sur suspenso-}
\l{\cel{3}risorti, tektoza, kun quar roti. - DFIS.}
\l{}
\l{\sou{karoto}. (\sou{bot.}) Planto ek la familio \textquotedbl{}umbeliferi\textquotedbl{}, di qua la}
\l{\cel{3}radiko (pulpoza) esas manjebla. - L.\cel{1}\sou{daucus carota}. -DEFI.}
\l{}
\l{\sou{karotido}. (\sou{anat.}) Singla ek la du grosa arteri qui portas la}
\l{\cel{3}sango a la kapo. - DEFIS.}
\l{}
\l{\sou{karotino}. (\sou{kemio)}. Hido-karburo, fuzebla ye 118o C., ye qua}
\l{\cel{3}konsistas la materio flava di la karoti, e qua absorbas}
\l{\cel{3}la oxo di aero.}
\l{}
\l{\sou{karpo}. I.\cel{1}\sou{(zool}.) Fisho qua vivas en la aquo sen-sala, ek la}
\l{\cel{3}genero \textquotedbl{}ciprini\textquotedbl{}. L.\cel{1}\sou{cyprinus carpio}. II. (\sou{anat.)}\cel{1}Parto di}
\l{\cel{3}la membro, supra od avana, inter la manuo e la avan-brakio.}
\l{\cel{3}- DEF.}
\l{}
\l{\sou{karpentar}. (\sou{trans.}) Facar la asemblajo ek ligno-pecegi, ek tra-}
\l{\cel{3}bi, qui uzesas en domo-konstrukturo por subtenar lua diversa}
\l{\cel{3}parti. - EFIS.}
\l{}
\l{karpino. (\sou{bot.}) Granda arboro qua kreskas en la foresti, e di}
\l{\cel{3}qua la ligno, harda, kompakta e blanka, demandesas por la}
\l{\cel{3}facaji da la tornisti, karpentisti, veturifisti, e c.}
\l{\cel{3}- L.\cel{1}\sou{carpinus betulus}. - ISL.}
\l{}
\l{karpologio. (\sou{bot.)}\cel{2}Parto di botaniko, qua havas kom studiajo,}
\l{\cel{3}la frukti.}
\l{}
\l{\sou{karto}. Quadrato o rektangulo ek sorto di papero rezistiva, ma}
\l{\cel{3}flexebla, facita ek plura papero-folii interglutinita.-DEFIRS.}
\l{}
\l{kartamo.\cel{2}(\sou{bot.}) Planto ek la familio \textquotedbl{}kompozaji\textquotedbl{}, di qua la}
\l{\cel{3}flori esas bele safrano-reda, uzata por tintar. - L. car-}
\l{\cel{3}thamus tinctorius. - FISL.}
\l{}
\l{kartamino. (\sou{kemio).}\cel{1}Koloro-materio di la flori, genero \textquotedbl{}kar-}
\l{\cel{3}tamo\textquotedbl{}. - DE.}
}
